% Ad Atticum 11.16 --- headnote
% ad-atticum-11-16, 3 June 47 BC, Brundisium

\textit{Cicero to Atticus, written from Brundisium on
the third day before the Nones of June 47 BC ---
3 June. The Perseus dateline header is corrupt
(\textit{Scr.\ Brundisi.\ iii Nou.\ luot.\ a.\ 707
(47)}), and the \textit{works.yaml} entry carries a
year-precision placeholder of \textit{-0047-01-25};
but the letter's own colophon at the close of \S5
preserves a clean \textit{iii Non.\ Iun.}, and the
content --- Caesar still detained at Alexandria,
peace still unsealed, the Pompeians of Asia and
Achaia now angling for Fufius's pardon --- fits the
beginning of June 47 BC well. The year-placeholder
in the manifest should be corrected to a day-precision
of -0047-06-03.}

\textit{The letter opens with cold reading of an
official document. Atticus has sent on a letter
purporting to come from Caesar; Cicero finds it
thinly written and full of marks against its
authenticity --- marks Atticus has noticed too. On
the question of going to meet Caesar he will follow
Atticus's advice, but he believes none of the
sustaining rumours: there is no real report of
Caesar's return, the men from Asia hear nothing of
peace, and the blows have fallen one after another
--- in Asia, in Illyricum, in the Cassius affair,
at Alexandria itself, at Rome, and in Italy. Even
if Caesar does come back, Cicero thinks the war
will be over before he arrives. The middle of the
letter is the same self-accusation as in 11.15,
now sharpened: the Achaian Pompeians who had once
shared his fear and his resolution are now suing
Fufius for pardon, and so the delay at Alexandria
has set their cause right and undone his. The
closing sections turn to private business ---
Quintus's renewed hostility, a stray sentence from
some friend at Patrae (\textit{Patris}) about
being not unwilling to be there given everything,
news passed via the freedman Cephalio who has
been held up for months by the weather; and a
delicate request, last of all, that Atticus and
Camillus together advise Terentia on her will,
since Philotimus has reported (\textit{vix
credibile}) that she is acting wickedly in matters
of debt. Cicero asks Atticus to write back even if
he can think of nothing: that very silence will
count for him as despair confirmed.}
